% p351_1; source PDF page 351; manually verified editable table.
\begin{tikzpicture}[x=1cm,y=-1cm,inner sep=0pt,outer sep=0pt]
\path[use as bounding box] (0,0) rectangle (16.000,11.725);
\draw[line width=.4pt] (0.000,0.000) rectangle (2.250,0.865);
\node[align=center,text width=1.990cm,font={\linespread{1}\fontsize{10}{13.50}\selectfont\bfseries}] at (1.125,0.432) {保险诈骗罪};
\draw[line width=.4pt] (2.250,0.000) rectangle (6.450,0.865);
\node[align=center,text width=3.940cm,font={\linespread{1}\fontsize{10}{13.50}\selectfont\bfseries}] at (4.350,0.432) {内容};
\draw[line width=.4pt] (6.450,0.000) rectangle (16.000,0.865);
\node[align=center,text width=9.290cm,font={\linespread{1}\fontsize{10}{13.50}\selectfont\bfseries}] at (11.225,0.432) {范例};
\draw[line width=.4pt] (0.000,0.865) rectangle (2.250,8.695);
\node[align=center,text width=1.990cm,font={\linespread{1}\fontsize{10}{13.50}\selectfont}] at (1.125,4.780) {行为主体};
\draw[line width=.4pt] (2.250,0.865) rectangle (6.450,6.880);
\node[align=left,text width=3.940cm,font={\linespread{1}\fontsize{10}{13.50}\selectfont}] at (4.350,3.873) {是特殊主体，限于投保人、被保险人、受益人。不具有主体身份的人骗取保险金的，可构成普通的诈骗罪};
\draw[line width=.4pt] (6.450,0.865) rectangle (16.000,6.880);
\node[align=left,text width=9.290cm,font={\linespread{1}\fontsize{10}{13.50}\selectfont}] at (11.225,3.873) {例1，个体户甲开办的汽车修理厂系某保险公司指定的汽车修理厂家。甲在为他人修理汽车时，多次夸大汽车毁损程度，向保险公司多报汽车修理费用，从保险公司骗取12万余元。对甲的行为应以诈骗罪论处。（2004年第5题）\\{}例2，甲公司给经理王某配了一辆轿车，轿车登记在公司名下，由王某使用。王某将车私下变卖，对甲公司领导谎称被盗，让甲公司去申请理赔。甲公司不知情，到保险公司理赔，获得5万元。王某构成职务侵占罪。王某对保险公司不构成保险诈骗罪的间接正犯，因为王某不具有投保人的身份，但可以构成诈骗罪的间接正犯};
\draw[line width=.4pt] (2.250,6.880) rectangle (6.450,8.695);
\node[align=center,text width=3.940cm,font={\linespread{1}\fontsize{10}{13.50}\selectfont}] at (4.350,7.787) {包括单位};
\draw[line width=.4pt] (6.450,6.880) rectangle (16.000,8.695);
\node[align=left,text width=9.290cm,font={\linespread{1}\fontsize{10}{13.50}\selectfont}] at (11.225,7.787) {某公司决定诈骗保险金，故意放火烧毁公司的仓库，并危害了公共安全，然后骗取保险金。对该公司定保险诈骗罪。由于单位不能构成放火罪，对公司直接责任人员定放火罪};
\draw[line width=.4pt] (0.000,8.695) rectangle (2.250,10.210);
\node[align=center,text width=1.990cm,font={\linespread{1}\fontsize{10}{13.50}\selectfont}] at (1.125,9.453) {侵害对象};
\draw[line width=.4pt] (2.250,8.695) rectangle (6.450,10.210);
\node[align=left,text width=3.940cm,font={\linespread{1}\fontsize{10}{13.50}\selectfont}] at (4.350,9.453) {商业保险。不包括社会保险};
\draw[line width=.4pt] (6.450,8.695) rectangle (16.000,10.210);
\node[align=left,text width=9.290cm,font={\linespread{1}\fontsize{10}{13.50}\selectfont}] at (11.225,9.453) {甲拾得乙的社保卡。甲在医院收费处交费时，冒用乙的社保卡。甲构成诈骗罪，不构成保险诈骗罪};
\draw[line width=.4pt] (0.000,10.210) rectangle (2.250,11.725);
\node[align=center,text width=1.990cm,font={\linespread{1}\fontsize{10}{13.50}\selectfont}] at (1.125,10.968) {着手的标准};
\draw[line width=.4pt] (2.250,10.210) rectangle (16.000,11.725);
\node[align=left,text width=13.490cm,font={\linespread{1}\fontsize{10}{13.50}\selectfont}] at (9.125,10.968) {虚构保险标的、制造保险事故，属于预备行为；只有到保险公司理赔，才是着手，才是实行行为};
\end{tikzpicture}
